% =====================================================================
%  CHAPTER 10: 儿童少年传染病
% =====================================================================
\chapter{儿童少年传染病}
\setcounter{chapter}{10}
\chapterintro{
掌握儿童少年常见传染病的流行病学特征；掌握学校传染病防控的基本原则和措施；掌握计划免疫在儿童传染病预防中的核心作用；熟悉新发传染病和再发传染病对儿童健康的威胁；了解学校传染病暴发的应急处理流程。
}

% ----- 10.1 概述 -----
\section{儿童少年传染病概述}

\begin{defbox}
\textbf{儿童少年传染病}：由各种病原体（病毒、细菌、真菌、寄生虫等）引起的能在人与人之间或人与动物之间传播的感染性疾病，是儿童少年群体的重要健康威胁。

\textbf{历史演变：}
\begin{itemize}[leftmargin=*]
    \item 20世纪初：传染病是儿童主要致病和死亡原因（天花、霍乱、麻疹、脊髓灰质炎等）
    \item 20世纪中期至今：计划免疫和公共卫生措施推广，许多传染病得到有效控制或消灭（如1980年全球消灭天花）
    \item 21世纪的新挑战：新发传染病（SARS、COVID-19等）和再发传染病（结核、水痘等）仍对儿童构成威胁
\end{itemize}
\end{defbox}

\subsection{儿童传染病高发的原因}

\begin{keybox}
\textbf{儿童是传染病的易感人群，原因包括：}
\begin{enumerate}[leftmargin=*]
    \item \imp{免疫系统发育不完善}——特异性免疫和非特异性免疫功能均未成熟
    \item \imp{学校集体生活环境}——人群密集、接触密切，有利于病原体传播
    \item \imp{卫生习惯尚未养成}——手卫生意识差、共用餐具等
    \item \imp{免疫空白}——部分疫苗未及时接种或抗体水平下降
    \item \imp{气候变化和人口流动}——增加了传染病传播风险
\end{enumerate}
\end{keybox}

% ----- 10.2 常见传染病 -----
\section{儿童少年常见传染病}

\begin{table}[H]
\centering
\caption{儿童少年常见传染病及特点}
\begin{tabularx}{\textwidth}{l>{\raggedright\arraybackslash}p{2.5cm}>{\raggedright\arraybackslash}X>{\raggedright\arraybackslash}p{2cm}}
\hline
\textbf{疾病} & \textbf{病原体} & \textbf{传播途径} & \textbf{主要特点} \\
\hline
手足口病 & 肠道病毒（EV71、CoxA16等）& 粪-口、接触、飞沫 & 发热、手足口部皮疹；5岁以下高发 \\
水痘 & 水痘-带状疱疹病毒 & 空气飞沫、接触 & 分批出现的皮疹（红斑疹→丘疹→水疱→结痂）；高度传染性 \\
流行性感冒 & 流感病毒 & 空气飞沫 & 发热、全身症状重；易发生暴发 \\
麻疹 & 麻疹病毒 & 空气飞沫（高度传染）& 发热、卡他症状、Koplik斑、皮疹 \\
流行性腮腺炎 & 腮腺炎病毒 & 空气飞沫 & 腮腺肿痛；可并发睾丸炎/卵巢炎 \\
肺结核 & 结核分枝杆菌 & 空气飞沫 & 咳嗽咳痰$>$2周、低热盗汗；卡介苗预防 \\
诺如病毒感染 & 诺如病毒 & 粪-口 & 呕吐、腹泻；冬季高发；极易在集体单位暴发 \\
\hline
\end{tabularx}
\end{table}

% ----- 10.3 传染病防控 -----
\section{儿童少年传染病的预防控制}

\subsection{计划免疫}

\begin{keybox}
\textbf{计划免疫（Planned Immunization）}：根据疫情监测和人群免疫状况分析，按照规定免疫程序，有计划地利用疫苗进行预防接种，以提高人群免疫水平、控制乃至消灭相应传染病。

\textbf{我国儿童计划免疫的里程碑成就：}
\begin{itemize}[leftmargin=*]
    \item 1978年开始实施计划免疫
    \item 2000年实现无脊髓灰质炎（维持至今）
    \item 2007年扩大免疫规划至14种疫苗预防15种疾病
    \item 5岁以下儿童乙肝表面抗原携带率从1992年的9.7\%降至2020年的$<$1\%
\end{itemize}

\textbf{接种疫苗是预防传染病\imp{最经济、最有效}的手段。}

\textbf{我国国家免疫规划疫苗程序（核心内容）：}
\begin{table}[H]
\centering
\caption{国家免疫规划主要疫苗及其预防疾病}
\begin{tabularx}{\textwidth}{l>{\raggedright\arraybackslash}p{3cm}>{\raggedright\arraybackslash}X}
\hline
\textbf{疫苗} & \textbf{接种月龄/年龄} & \textbf{预防疾病} \\
\hline
卡介苗（BCG） & 出生时 & 结核病（特别是粟粒性结核和结核性脑膜炎） \\
乙肝疫苗（HepB） & 0、1、6月龄 & 乙型病毒性肝炎 \\
脊灰灭活疫苗（IPV）/减毒活疫苗（bOPV）& 2、3、4月龄 + 4岁 & 脊髓灰质炎 \\
百白破疫苗（DTaP） & 3、4、5月龄 + 18月龄 & 百日咳、白喉、破伤风 \\
麻腮风疫苗（MMR） & 8月龄（麻风）+ 18月龄（麻腮风）& 麻疹、流行性腮腺炎、风疹 \\
乙脑减毒活疫苗（JEV-L）& 8月龄、2岁 & 流行性乙型脑炎 \\
A群流脑多糖疫苗 & 6、9月龄 & 流行性脑脊髓膜炎 \\
A+C群流脑多糖疫苗 & 3岁、6岁 & 流行性脑脊髓膜炎 \\
甲肝减毒活疫苗 & 18月龄 & 甲型病毒性肝炎 \\
白破疫苗（DT） & 6岁（加强） & 白喉、破伤风 \\
\hline
\end{tabularx}
\end{table}

\textbf{入学查验接种证：}小学入学时必须查验预防接种证，发现漏种者应及时补种。这是学校传染病防控的第一道关口。
\end{keybox}

\subsection{学校传染病防控}

\begin{keybox}
\textbf{学校传染病防控的关键措施：}

\begin{enumerate}[leftmargin=*]
    \item \imp{入学查验预防接种证}——确保入学前完成基础免疫，发现漏种及时补种

    \item \imp{晨检与因病缺课追踪}——早期发现疑似传染病病例
    \begin{itemize}
        \item 晨检：每天早晨对学生进行健康检查（测体温、观察症状）
        \item 因病缺课追踪：追踪缺课学生病因，发现传染病线索
    \end{itemize}

    \item \imp{疫情报告}——发现法定传染病应立即向当地疾病预防控制中心和上级教育行政部门报告

    \item \imp{隔离与消毒}
    \begin{itemize}
        \item 传染源管理：确诊或疑似患者及时隔离、治疗
        \item 切断传播途径：教室通风换气、消毒、手卫生
        \item 保护易感者：应急接种、健康教育
    \end{itemize}

    \item \imp{健康教育与健康促进}
    \begin{itemize}
        \item 良好个人卫生习惯的培养（勤洗手、咳嗽礼仪等）
        \item 增强免疫力（均衡营养、充足睡眠、适度锻炼）
    \end{itemize}
\end{enumerate}
\end{keybox}

\subsection{新发传染病对儿童的挑战}

\begin{defbox}
\textbf{新发传染病（Emerging Infectious Diseases）}：近年来新出现或重新出现且对人类健康构成威胁的传染病，如SARS（2003）、甲型H1N1流感（2009）、COVID-19（2019）等。

\textbf{对儿童的影响：}
\begin{itemize}[leftmargin=*]
    \item 学校停课影响儿童学习和心理健康
    \item 儿童感染后可能表现为轻症但成为传播链条的重要环节
    \item 疫苗接种中断可能导致其他传染病暴发
\end{itemize}
\end{defbox}

% ----- 复习题 -----
\section{复习题}

\begin{enumerate}[leftmargin=*, itemsep=8pt]
    \item \textbf{【名词解释】}计划免疫；晨检；新发传染病
    \item \textbf{【简答题】}简述国家免疫规划的主要内容及入学查验接种证的意义。
    \item \textbf{【简答题】}简述学校传染病防控的关键措施。
    \item \textbf{【简答题】}简述儿童易患传染病的主要原因。
    \item \textbf{【简答题】}简述计划免疫在儿童传染病防控中的意义。
\end{enumerate}

\subsection*{参考答案要点}

\begin{enumerate}[leftmargin=*, itemsep=5pt]
    \item \textbf{计划免疫}：按规定免疫程序有计划地利用疫苗进行预防接种，提高人群免疫水平、控制乃至消灭传染病。\textbf{晨检}：每天早晨对学生进行健康检查，早期发现疑似传染病病例。\textbf{新发传染病}：新出现或重新出现对人类健康构成威胁的传染病，如SARS、COVID-19等。

    \item 国家免疫规划包括：卡介苗（结核病）、乙肝疫苗（乙肝）、脊灰疫苗（脊灰）、百白破疫苗（百日咳、白喉、破伤风）、麻腮风疫苗（麻疹、腮腺炎、风疹）、乙脑疫苗（乙脑）、流脑疫苗（流脑）、甲肝疫苗（甲肝），在出生至6岁间按程序接种。入学查验接种证确保入学前完成基础免疫，是学校防控第一道关口。

    \item (1)入学查验接种证；(2)晨检与因病缺课追踪（早发现）；(3)疫情报告（法定传染病应报告）；(4)隔离与消毒（传染源管理、切断传播途径、保护易感者）；(5)健康教育（洗手、咳嗽礼仪等）。

    \item (1)免疫系统发育不完善；(2)学校集体生活（人群密集、接触密切）；(3)卫生习惯未养成；(4)免疫空白（漏种或抗体水平下降）；(5)气候变化和人口流动增加传播风险。

    \item 计划免疫是最经济最有效的传染病预防手段：已实现全球消灭天花、中国无脊灰、5岁以下儿童乙肝携带率从9.7\%降至$<$1\%。建议所有适龄儿童按时完成国家免疫规划疫苗接种。
\end{enumerate}

\newpage
